% 拉普拉斯-贝尔特拉米算子（综述）
% license CCBYSA3
% type Wiki

本文根据 CC-BY-SA 协议转载翻译自维基百科\href{https://en.wikipedia.org/wiki/Laplace\%E2\%80\%93Beltrami_operator}{相关文章}。

在微分几何中，拉普拉斯–贝尔特拉米算子（Laplace–Beltrami operator）是对拉普拉斯算子（Laplace operator）的一种推广，它作用于定义在欧几里得空间子流形上的函数，更一般地，也可定义在黎曼流形（Riemannian manifold）及伪黎曼流形（pseudo-Riemannian manifold）上。该算子以皮埃尔–西蒙·拉普拉斯（Pierre-Simon Laplace）与欧金尼奥·贝尔特拉米（Eugenio Beltrami）的名字命名。

对于任意定义在欧几里得空间$\mathbb{R}^{n}$上的二次可微实值函数$f$，拉普拉斯算子（又称拉普拉斯算符，Laplacian）将$f$映射为其梯度向量场的散度，即$f$关于$\mathbb{R}^{n}$的一个正交基的$n$个纯二阶偏导数之和。

与拉普拉斯算子类似，拉普拉斯–贝尔特拉米算子也定义为梯度的散度，它是一个线性算子，将函数映射为函数。

该算子可推广至张量场的情形：定义为协变导数（covariant derivative）的散度；或者，也可推广为作用在微分形式上的算子，使用散度与外微分（exterior derivative）定义。所得到的算子称为拉普拉斯–德拉姆算子（Laplace–de Rham operator），以乔治·德拉姆（Georges de Rham）的名字命名。
\subsection{细节}
与普通拉普拉斯算子类似，拉普拉斯–贝尔特拉米算子（Laplace–Beltrami operator）定义为（黎曼意义下的）梯度的散度：
\[
\Delta f = \mathrm{div}(\nabla f)~
\]
在局部坐标下，可以给出其显式形式。

首先，设$M$是一个定向黎曼流形（oriented Riemannian manifold）。
定向允许我们在$M$上指定一个确定的体积形式（volume form），在定向坐标系$\{x^{i}\}$下表示为：
\[
\operatorname{vol}_{n}
:= \sqrt{|g|}\; dx^{1} \wedge \cdots \wedge dx^{n},~
\]
其中$|g| := |\det(g_{ij})|$为度量张量（metric tensor）的行列式的绝对值，
$dx^{i}$为与切丛基
\[
\partial_{i} := \frac{\partial}{\partial x^{i}}~
\]
对偶的 1-形式，符号$\wedge$表示外积（wedge product）。

向量场$\mathbf{X}$在流形上的散度（divergence）定义为满足以下条件的标量函数 $\nabla \cdot \mathbf{X}$：
\[
(\nabla \cdot \mathbf{X}) \, \operatorname{vol}_{n}
:= L_{\mathbf{X}} (\operatorname{vol}_{n}),~
\]
其中$L_{\mathbf{X}}$表示沿向量场$\mathbf{X}$的李导数（Lie derivative）。

在局部坐标中，可得：
\[
\nabla \cdot \mathbf{X}
= \frac{1}{\sqrt{|g|}} \, \partial_{i}
  \left( \sqrt{|g|} \, X^{i} \right),~
\]
其中及以下均采用爱因斯坦求和约定（Einstein summation convention），即对重复指标$i$进行求和。

标量函数$f$的梯度（gradient）是一个向量场$\operatorname{grad} f$，它通过流形上的内积$\langle \cdot, \cdot \rangle$定义为：
\[
\langle \operatorname{grad} f(x), v_{x} \rangle = df(x)(v_{x}),~
\]
其中$v_{x}$为锚定在点$x$的切向量（即$v_{x} \in T_{x}M$），$df$为函数$f$ 的外微分（exterior derivative），它是一个作用于$v_{x}$的 1-形式。

在局部坐标下，有：
\[
(\operatorname{grad} f)^{i}
= \partial^{i} f
= g^{ij} \partial_{j} f,~
\]
其中$g^{ij}$为度量张量$g_{ij}$的逆矩阵分量，满足$g^{ij} g_{jk} = \delta^{i}_{k}$，$\delta^{i}_{k}$为克罗内克 δ 符号（Kronecker delta）。

结合梯度与散度的定义，作用于标量函数$f$的拉普拉斯–贝尔特拉米算子（Laplace–Beltrami operator）在局部坐标下的表达式为：
\[
\Delta f
= \frac{1}{\sqrt{|g|}} \,
  \partial_{i}
  \left(
    \sqrt{|g|} \, g^{ij} \, \partial_{j} f
  \right)~
\]
若流形 $M$ 非定向（non-oriented），上述推导仍然完全成立，唯一的区别在于体积形式（volume form）需被体积元（volume element）所取代，即此时用密（density）而非微分形式来定义积分测度。由于梯度与散度本身并不依赖于定向的选择，因此拉普拉斯–贝尔特拉米算子$\Delta$也独立于这种附加结构。
\subsection{形式自伴性}
外微分算子$d$ 与 $-\nabla$在形式上互为伴随（formal adjoints），其意义在于：对于任意紧支函数（compactly supported function）$f$，
有
\[
\int_{M} df(X) \, \operatorname{vol}_{n}
= -\int_{M} f \, \nabla \cdot X \, \operatorname{vol}_{n},~
\]
其中最后一个等式可由\emph{斯托克斯定理}（Stokes' theorem）直接得出。

通过对偶化（dualization），可得：
\[
\int_{M} f \, \Delta h \, \operatorname{vol}_{n}
= -\int_{M} \langle df, dh \rangle \, \operatorname{vol}_{n},
\tag{2}~
\]
其中$\langle df, dh \rangle$表示 1-形式$df$与$dh$在度量诱导下的内积。该式对所有紧支函数$f$与$h$成立。反之，式（2）可以完全刻画拉普拉斯–贝尔特拉米算子，
即它是唯一满足该性质的算子。

由此可知，拉普拉斯–贝尔特拉米算子是负定的（negative）且形式自伴的（formally self-adjoint），即对于所有紧支函数$f$与$h$，有：
\[
\int_{M} f \, \Delta h \, \operatorname{vol}_{n}
= -\int_{M} \langle df, dh \rangle \, \operatorname{vol}_{n}
= \int_{M} h \, \Delta f \, \operatorname{vol}_{n}.~
\]
由于以上定义下的拉普拉斯–贝尔特拉米算子是负的而非正的，因此在许多文献与物理学中，常采用相反符号的定义，以使算子成为正定形式。
\subsection{拉普拉斯–贝尔特拉米算子的特征值}
设$M$为一个无边界的紧黎曼流形（compact Riemannian manifold without boundary）。考虑以下特征值方程：
\[
-\Delta u = \lambda u~
\]
其中$u$为对应于特征值$\lambda$的特征函数（eigenfunction）。

由前面所证的形式自伴性（formal self-adjointness）可知，特征值$\lambda$必为实数。由于流形$M$的紧性（compactness），可进一步证明特征值是离散的（discrete），并且每个特征值对应的特征函数空间（eigenspace）是有限维的。

取常数函数为特征函数可得$\lambda = 0$是一个特征值。此外，由于我们考虑的是$-\Delta$，经分部积分（integration by parts）可得
\[
\lambda \ge 0~
\]
更具体地，若我们将特征值方程两边同乘以特征函数$u$，并在$M$上积分（记$dV = \operatorname{vol}_{n}$），则：
\[
-\int_{M} (\Delta u)\, u\, dV = \lambda \int_{M} u^{2}\, dV~
\]
对左边的项应用散度定理（divergence theorem），并注意到$M$无边界，得：
\[
-\int_{M} (\Delta u)\, u\, dV = \int_{M} |\nabla u|^{2}\, dV~
\]
结合以上两式，可得：
\[
\int_{M} |\nabla u|^{2}\, dV = \lambda \int_{M} u^{2}\, dV~
\]
从而结论 $\lambda \ge 0$ 立即成立。

\textbf{Lichnerowicz 定理.}

一个由 André Lichnerowicz 证明的重要结果指出：设$M$为一个$n$维（$n \ge 2$）紧黎曼流形，其Ricci 曲率（Ricci curvature）满足下界：
\[
\operatorname{Ric}(X, X) \ge \kappa\, g(X, X), \quad \kappa > 0,~
\]
其中$g(\cdot, \cdot)$为度量张量（metric tensor），$X$为任意切向量。则第一个正特征值$\lambda_{1}$满足不等式：
\[
\lambda_{1} \ge \frac{n}{n-1}\, \kappa~
\]
该下界是锐的（sharp），并在球面$\mathbb{S}^{n}$上取到。

事实上，在二维情形$\mathbb{S}^{2}$ 上，$\lambda_{1}$的特征子空间是三维的，由坐标函数$x_{1}, x_{2}, x_{3}$（即$\mathbb{R}^{3}$的坐标函数限制到 $\mathbb{S}^{2}$上）所张成。在球坐标系$(\theta, \phi)$下，令
\[
x_{3} = \cos \phi = u_{1}~
\]
由球面拉普拉斯算子的公式可直接计算得：
\[
-\Delta_{\mathbb{S}^{2}} u_{1} = 2u_{1}~
\]
这表明 Lichnerowicz 定理的下界在二维情形下确实可以达到。

\textbf{Obata 定理.}

反之，Morio Obata 证明了如下结论：若$M$是$n$维紧黎曼流形且无边界，并且其第一个正特征值满足：
\[
\lambda_{1} = \frac{n}{n-1}\, \kappa~
\]
则$M$与半径为
\[
\sqrt{\frac{n-1}{\kappa}}~
\]
的$n$维球面$\mathbb{S}^{n}\!\left(\sqrt{\frac{n-1}{\kappa}}\right)$ 同构等距（isometric）。

这些命题的证明可见于 Isaac Chavel 的著作。类似的锐界（sharp bounds）也存在于其他几何情形中，例如紧致 CR 流形上的Kohn 拉普拉斯算子（Kohn Laplacian, Joseph J. Kohn），其应用包括此类流形在复空间$\mathbb{C}^{n}$中的整体嵌入问题。
